% ============ 图 1　工具调用控制系统的结构示意图 ============
%
% 网格（单位 cm）：列 x = 0 / 3.6 / 7.2，行 y = 0 / -2.0 / -4.0
% 模块框 2.4 × 0.95，外框 100 从 (-1.6,0.95) 到 (9.7,-4.75)
% 要挪模块只改下面九个 \node 的坐标，连线用的是锚点会自动跟随。
%
%          左列      中列      右列
%   上行   101       105       102
%   中行   109       104       103
%   下行   108       107       106
%
% 全图只有一处连线交叉：107→101 的竖段与 108→104 的横段交于 (1.8,-2.9)。
% 这是布局结构决定的，按普通交叉绘制，不加跨线弧。
% ==============================================================
\begin{tikzpicture}[x=1cm,y=1cm]

  % ---- 外框 100 ----
  \draw[pline] (-1.6,0.95) rectangle (9.7,-4.75);
  \node[anchor=south west, font=\normalsize] at (-1.58,1.0) {100};

  % ---- 智能体运行时内部的九个单元 ----
  \node[pmod] (n101) at (0   , 0   ) {101};
  \node[pmod] (n105) at (3.6 , 0   ) {105};
  \node[pmod] (n102) at (7.2 , 0   ) {102};
  \node[pmod] (n109) at (0   ,-2.0 ) {109};
  \node[pmod] (n104) at (3.6 ,-2.0 ) {104};
  \node[pmod] (n103) at (7.2 ,-2.0 ) {103};
  \node[pmod] (n108) at (0   ,-4.0 ) {108};
  \node[pmod] (n107) at (3.6 ,-4.0 ) {107};
  \node[pmod] (n106) at (7.2 ,-4.0 ) {106};

  % ---- 运行时之外：大语言模型与工具集合 ----
  \node[pmodsmall] (n110) at (-3.6,0) {110};
  \node[pmodsmall] (n111) at (10.8,0) {111};

  % ---- 连线 ----
  \draw[plr] (n110.east)  -- (n101.west);              % 110 ↔ 101
  \draw[pl]  (n101.east)  -- (n105.west);              % 101 → 105
  \draw[pl]  (n105.east)  -- (n102.west);              % 105 → 102
  \draw[plr] (n102.east)  -- (n111.west);              % 102 ↔ 111
  \draw[pl]  (n102.south) -- (n103.north);             % 102 → 103
  \draw[pl]  (n103.west)  -- (n104.east);              % 103 → 104
  \draw[pl]  (n104.north) -- (n105.south);             % 104 → 105
  \draw[pl]  (n104.west)  -- (n109.east);              % 104 → 109
  \draw[pl]  (n106.west)  -- (n107.east);              % 106 → 107
  \draw[plr] (n107.west)  -- (n108.east);              % 107 ↔ 108

  % 102 → 106：自右侧 8.9 的通道绕过 103
  \draw[pl] ([yshift=2mm]n102.east) -- (8.9,0.2) -- (8.9,-4.0) -- (n106.east);

  % 107 → 101：走左列与中列之间 1.8 的通道
  \draw[pl] ([yshift=2mm]n107.west) -- (1.8,-3.8) -- (1.8,-0.2)
            -- ([yshift=-2mm]n101.east);

  % 108 → 104：自 108 上沿经 -2.9 的通道折向 104 下沿
  \draw[pl] (n108.north) -- (0,-2.9) -- (3.6,-2.9) -- (n104.south);

\end{tikzpicture}
